% !TeX root = booklet-03.tex
\ifCompleteNotes
\section*{راهنمای بخش سوم، مشخصه انتقال ولتاژ}
\addcontentsline{toc}{section}{راهنمای بخش سوم، مشخصه انتقال ولتاژ}
\else
\chapter*{راهنمای مطالعه}
\addcontentsline{toc}{chapter}{راهنمای مطالعه}
\fi
این جزوه، ارتباط ولتاژ ورودی و خروجی یک مدار منطقی را با جریان
ترانزیستورهای آن توضیح می‌دهد. آشنایی قبلی با \lr{Digital Electronics}
فرض نشده است. معنای صفر و یک، عملکرد \lr{Inverter}، شرط تعادل جریان،
نمودار تابع و محاسبات درجه دوم پیش از تحلیل مرور می‌شوند.

این جزوه، مقدمه و تمام بخش \secnum{۲.۵} کتاب را بر اساس
صفحات چاپی ۸۷ تا آغاز ۹۳ و اسلایدهای ۳ تا ۱۴
\lr{CMOS Inverter DC Characteristics} پوشش می‌دهد \cite{weste,dcslides}. استخراج روابط انتقال،
انتخاب ریشه‌ها، کنترل ناحیه کاری و مثال‌های عددی، توضیح آموزشی
روابط مدل کتاب‌اند. همه پارامترهای مثال‌های افزوده صریحاً فرضی
معرفی می‌شوند. نمودارهای محاسبه‌شده، نتیجه همین مدل‌ها هستند.

ابتدا مرور پیش‌نیازها و سپس ترتیب ولتاژها، جریان‌ها، ناحیه کاری و
حل معادله را دنبال کنید. مشخصه ایده‌آل در نقطه میانی یک ظرافت
مهم دارد. برابری دو جریان اشباع، ورودی را تعیین می‌کند و خروجی
در بازه‌ای از ولتاژها می‌تواند تعادل داشته باشد. این بازه استخراج
می‌شود و اثر \lr{Channel Length Modulation (CLM)} بر شیب نیز
به‌طور جداگانه توضیح داده می‌شود.

سپس نسبت \(\beta\)، جابه‌جایی آستانه، تفاوت تعریف نسبت قدرت در
دو مدل جریان و تحلیل \lr{Equivalent Inverter} بررسی می‌شوند.
بحث \lr{Noise Margin}، سطوح منطقی، استخراج نقاط شیب منفی یک و
حل کامل تمرین \secnum{۲.۱۷} نیز آمده‌اند. تفاوت تحلیل نویز ایستا و زمانی
با اشاره به اسلاید ۱۱ توضیح داده می‌شود.
بخش \lr{Pass Transistor}، حد انتقال ولتاژ، اثر بدنه،
\lr{Transmission Gate} و حل تمرین \secnum{۲.۲۰} را توضیح می‌دهد.
تمرین \secnum{۲.۱۶} در استخراج تحلیلی ناحیه‌های \lr{B} و \lr{D}
حل شده و در مرور پایانی مشخص شده است. نقشه پوشش منابع و هر سه
تمرین سری دوم فصل دوم نیز در پایان آمده است.
شماره روابط و شکل‌های این جزوه مستقل از کتاب است.

